% ============================================================================
\clearpage
\section*{v4.9.0-r42 Source-Only: A Compensator of the Millennium Matrix as Chirality Projector}
\addcontentsline{toc}{section}{v4.9.0-r42 Source-Only: A Compensator of the Millennium Matrix as Chirality Projector}
\label{sec:v490-r42-mm-compensator-chirality}
% ============================================================================

\paragraph{Purpose.}
Layer v4.9.0-r41 placed the weak \(SU(2)\) on the right of \(\HH\) and obtained only weak doublets.  The author states (2026-10-23) that
the Millennium matrix of \S\ref{sec:millennium-hc}, which was to carry several Heisenberg compensators, was meant for the missing half.
This layer records the pre-registered test H27 (note \texttt{llm/Vorab\_Notiz\_H27\_MM\_Kompensator\_Chiralitaet.md}, commit
\texttt{614ac3e}, fixed before computing).  The charge pattern was seen analytically before the note and is disclosed there; the
computation checks it and is not a prediction.  Earlier layers are not rewritten.

\noindent\textbf{Source-only status.}
This v4.9.0-r42 layer is source-only.  It does not build a PDF, does not assign a new DOI,
and does not authorize a public release.  The v4.9.0 Version DOI
\href{https://doi.org/10.5281/zenodo.23534954}{10.5281/zenodo.23534954} and the Concept DOI
\href{https://doi.org/10.5281/zenodo.19210728}{10.5281/zenodo.19210728} are unchanged by it.

\subsection*{v4.9.0-r42 construction}
\label{subsec:v490-r42-construction}

On \(\mathbb C\otimes\HH\otimes\OO\cong\mathbb C^{32}\) with the distribution D of v4.9.0-r41, put
\[
C_{\rm ch}=\tfrac12(\mathrm{id}+J),\qquad J=i\,L_{e_7},\qquad T_a'=C_{\rm ch}T_a,
\]
\[
\tfrac{Y}{2}=Q-\tfrac12+(1-C_{\rm ch})T_3,\qquad Q_{\rm em}=T_3+Q-\tfrac12 ,
\]
with \(i\) the outer unit and \(L_{e_7}\) the inner unit, so \(J^2=1\) and \(C_{\rm ch}\) has the form of the Heisenberg compensator.

\subsection*{v4.9.0-r42 results}
\label{subsec:v490-r42-results}

\paragraph{Theoremlet v4.9.0-r42.1.}
\begin{enumerate}[leftmargin=2em]
\item[(a)] \(C_{\rm ch}\) is a Hermitian projector of rank \(16\); it commutes with colour, the Lorentz group, the right action of \(\HH\)
and \(Q\), and its image is the sum of the \(Q\)-sectors \(0\) and \(\tfrac23\).
\item[(b)] Under colour, the weak \(SU(2)'\) and \(Y/2\) the space splits into
\begin{center}
\small
\begin{tabular}{@{}llll@{}}
\toprule
\((SU(3),SU(2)',Y/2)\) & \(Q_{\rm em}\) & state & complex dimension\\
\midrule
\((3,2,\tfrac16)\) & \(\tfrac23,-\tfrac13\) & \((u,d)_L\) & \(12\)\\
\((1,2,-\tfrac12)\) & \(0,-1\) & \((\nu,e)_L\) & \(4\)\\
\((\bar3,1,-\tfrac23)\) & \(-\tfrac23\) & \(u^c\) & \(6\)\\
\((\bar3,1,\tfrac13)\) & \(\tfrac13\) & \(d^c\) & \(6\)\\
\((1,1,1)\) & \(1\) & \(e^c\) & \(2\)\\
\((1,1,0)\) & \(0\) & \(\nu^c\) & \(2\)\\
\bottomrule
\end{tabular}
\end{center}
that is, sixteen Weyl states with two Lorentz components each; \(Y/2\) commutes with \(SU(2)'\), colour and the Lorentz group, and
\(Q_{\rm em}=T_3'+Y/2\).
\item[(c)] Over the sixteen states, \(\operatorname{Tr}Y\), \(\operatorname{Tr}Y^3\) and the \(SU(2)^2U(1)\) and \(SU(3)^2U(1)\) traces vanish.
\item[(d)] \(SU(2)'\) acts only on the image of \(C_{\rm ch}\), and \(TC_{\rm ch}T^{-1}=1-C_{\rm ch}\).
\item[(e)] \(\mathcal M=\lambda_0C_{\rm ch}+\lambda_1(1-C_{\rm ch})\), \(\lambda_0\neq\lambda_1\), has \(C_{\rm ch}\) as spectral projector, as in
Theorem~\ref{thm:hc-spectral-projector}.
\end{enumerate}

\paragraph{Proof sketch.}
By explicit complex matrices (script \texttt{llm/k83\_test\_h27.py}, deviations at most \(10^{-15}\)).  \(C_{\rm ch}\) equals \(P_-\) of
v4.9.0-r35.1(c).  A first run reported (b) as refuted because the candidate values of \(Y/2\) were rounded to nine digits while the
null-space test required \(10^{-10}\), which lost all non-dyadic values; only this rounding was corrected, and both outputs are kept.
\(\square\)

\paragraph{Verdict and reading.}
H27a--e are confirmed, (b) after the correction of the script; (e) is a control.
\begin{enumerate}[leftmargin=2em]
\item A compensator of the Companion's form \(\tfrac12(\mathrm{id}+J)\), with \(J\) the product of the outer and the inner unit, completes
the doublets of v4.9.0-r41 to the full left-handed content of one generation with right-handed neutrino.  The weak \(SU(2)\) acts only
on its image: in the model, this is chirality.  A compensator of the Millennium matrix separates the chiralities.
\item The frame change \(T\) exchanges image and kernel of the compensator.  With v4.9.0-r41.1(c), matter frame and tachyon frame are in
the model particle and antiparticle with mirrored chirality.  This is a reading.
\item The Millennium matrix of \S\ref{sec:millennium-hc} is \(16\times16\), and one generation with \(\nu_R\) has sixteen Weyl states.  This is
recorded, not interpreted.
\item This is a construction under choices, not a derivation: \(J\) and the formula for \(Y/2\) are chosen so that the known result
appears, and the pattern was seen in advance.  It shows that the building blocks of the Companion (compensator, two units, two sides,
triality) suffice.  Anomaly cancellation follows from the content.  Similar constructions exist in the literature on
\(\mathbb C\otimes\HH\otimes\OO\).  Why this choice, symmetry breaking, masses and mixings are not addressed.
\end{enumerate}

\subsection*{v4.9.0-r42 status in the three frames}
\label{subsec:v490-r42-frames}

\begin{center}
\small
\begin{tabular}{@{}>{\raggedright\arraybackslash}p{0.30\textwidth}>{\raggedright\arraybackslash}p{0.62\textwidth}@{}}
\toprule
result & M / T / P\\
\midrule
r42.1(a)--(c) & stated in M\\
r42.1(d) & relates M and T: \(T\) exchanges the chiral halves; with r41.1(c) the T frame reads the conjugate states with mirrored chirality; P not evaluated\\
\bottomrule
\end{tabular}
\end{center}

\subsection*{v4.9.0-r42 register entries}
\label{subsec:v490-r42-register}

\begin{center}
\small
\begin{tabular}{@{}>{\raggedright\arraybackslash}p{0.34\textwidth}>{\raggedright\arraybackslash}p{0.58\textwidth}@{}}
\toprule
statement & status after v4.9.0-r42\\
\midrule
content of one generation & constructed under choices from compensator, two units, two sides (r42.1(b)); not derived\\
chirality & in the model: the weak \(SU(2)\) acts on the image of one compensator (r42.1(d))\\
Millennium matrix & a two-projector reading with \(C_{\rm ch}\) as spectral projector (r42.1(e)); the dimension \(16\) recorded\\
\bottomrule
\end{tabular}
\end{center}

\subsection*{v4.9.0-r42 ledger and guard update}
\label{subsec:v490-r42-ledger-guard}

\[
R_{\rm ch}^{v490\text{-}r42}=(1,0,0)
\]
for (content of one generation constructed; choice of \(J\) and \(Y\) not derived; breaking, masses and mixings open).

The layer spends no guard:
\[
R_{\rm glob}^{(v490\text{-}r42,+)}=(0,1,1,1,1),
\qquad
\Gamma_{\rm guard}^{v490\text{-}r42}=(1,1,1,1,1),
\qquad
\beta_4=1/256\ \text{unspent}.
\]
It does not imply exact finite \(C5\) zero closure, exact all-readout terminality,
\(C6/T2\), a global photon-clock/RH theorem, or global all-readout closure.
